\chapter{Conception fonctionnelle}
\label{chap:conception-fonctionnelle}

Ce chapitre a pour but de faire la modélisation du système AOS-ETL à partir de diagrammes UML. On y trouve un diagramme de cas d'utilisation précisant les acteurs principaux et leurs interactions avec le système, des diagrammes de séquence relatant les processus clés, un diagramme de classes destiné à représenter les entités fondamentales et leurs relations.

%% ============================================================
\section{Modélisation du système}
\label{sec:modelisation}
%% ============================================================

\subsection{Définition d'UML}

Le langage de modélisation unifié (UML) est un langage visuel standardisé lequel se prête à spécifier, visualiser, construire et documenter les systèmes logiciels, et permet de comprendre de façon claire la structure et le comportement d'un système, en diminuant les ambiguïtés et en garantissant la cohérence tout au long du projet~\cite{uml}.

\begin{figure}[H]
    \centering
    
\includegraphics[width=0.7\textwidth]{uml.jpeg}
    \caption{Les différents types de diagrammes UML}
    \label{fig:uml-types}
\end{figure}

\subsection{Diagramme de cas d'utilisation}

Le diagramme de cas d'utilisation ci-dessous présente les acteurs du système et leurs interactions principales avec le framework AOS-ETL.

\begin{figure}[H]
    \centering
    \includegraphics[width=\textwidth, trim=0 240 0 40, clip]{phase-mapping-use-case.pdf}
    \caption{Diagramme de cas d'utilisation — Phase Mapping}
    \label{fig:use-case-mapping}
\end{figure}

\begin{figure}[H]
    \centering
    \includegraphics[width=\textwidth, trim=0 400 0 40, clip]{post-mapping-use-case.pdf}
    \caption{Diagramme de cas d'utilisation — Phase Post-mapping}
    \label{fig:use-case-post-mapping}
\end{figure}

Les acteurs identifiés sont~:
\begin{itemize}
    \item \textbf{Responsable d'intégration~:} utilisateur principal qui upload les fichiers CSV, interagit avec l'agent de mapping, lance la génération des pipelines et l'exécution des imports ;
    \item \textbf{Agent IA (Claude)~:} acteur secondaire qui analyse les CSV, résout les dépendances FK, génère les manifestes et les pipelines.
\end{itemize}

\subsection{Raffinement des cas d'utilisation}

Le tableau~\ref{tab:raffinement} détaille les actions possibles par acteur.

\begin{table}[H]
  \centering
  \renewcommand{\arraystretch}{1.4}
  \begin{tabular}{p{3.5cm} p{4cm} p{6cm}}
    \toprule
    \rowcolor{gray!20}
    \textbf{Acteur} & \textbf{Action} & \textbf{Description} \\
    \midrule
    \rowcolor{gray!5}
    Responsable d'intégration & Démarrer une session & Choisit l'entité cible et configure la connexion BD \\
    Responsable d'intégration & Déposer des fichiers CSV & Upload un ou plusieurs fichiers dans la session \\
    \rowcolor{gray!5}
    Responsable d'intégration & Résoudre les FK & Répond aux questions de l'agent sur les ambiguïtés \\
    Responsable d'intégration & Consulter les fichiers & Vérifie les fichiers déposés dans la session \\
    \rowcolor{gray!5}
    Responsable d'intégration & Consulter le manifest & Revue du mapping proposé par l'agent \\
    Responsable d'intégration & Générer les pipelines & Lance la génération avec paramètres (paths, contexte) \\
    \rowcolor{gray!5}
    Responsable d'intégration & Exécuter le workflow & Lance l'import via hop-run.sh \\
    Responsable d'intégration & Consulter les résultats & Vérifie les lignes importées et les erreurs \\
    \midrule
    \rowcolor{gray!5}
    Agent IA (Claude) & Analyser le CSV & Profile les colonnes, détecte les types et FK \\
    Agent IA (Claude) & Résoudre les FK & Vérifie en BD l'existence des valeurs référencées \\
    \rowcolor{gray!5}
    Agent IA (Claude) & Produire le manifest & Écrit le fichier manifest.json \\
    Agent IA (Claude) & Générer les pipelines & Produit les fichiers .hpl/.hwf \\
    \bottomrule
  \end{tabular}
  \caption{Raffinement des cas d'utilisation}
  \label{tab:raffinement}
\end{table}

\subsection{Diagrammes de séquence}

\subsubsection{Séquence — Phase Mapping}

\begin{figure}[H]
    \makebox[\textwidth][c]{\includegraphics[width=1.25\textwidth]{uml-sequence-mapping.png}}
    \caption{Diagramme de séquence — session de mapping via l'agent IA}
    \label{fig:seq-mapping}
\end{figure}

\clearpage

\subsubsection{Séquence — Phase Post-mapping}

\begin{figure}[H]
    \makebox[\textwidth][c]{\includegraphics[width=1.25\textwidth]{uml-sequence-postmapping.png}}
    \caption{Diagramme de séquence — génération et exécution des pipelines}
    \label{fig:seq-postmapping}
\end{figure}

\clearpage

\subsection{Diagramme de classes}

\begin{figure}[H]
    \centering
    \includegraphics[width=0.9\textwidth]{uml-class-diagram.png}
    \caption{Diagramme de classes du framework AOS-ETL}
    \label{fig:class-diagram}
\end{figure}

%% ============================================================
\section*{Conclusion}
%% ============================================================

Ce chapitre a offert une vue fonctionnelle du système au moyen de diagrammes UML. Les cas d'utilisation ont permis de repérer les interactions entre le responsable d'intégration et le système, les diagrammes de séquence ont détaillé les flux principaux et le diagramme de classes a formalisé les entités majeures du framework. Le chapitre suivant explique l'architecture technique de la solution.
